\documentclass[11pt,oneside]{article}
% Standalone ProveIt research note.  The source intentionally has no external
% inputs so that the TeX/PDF pair is directly reproducible after extraction.
\usepackage[a4paper,margin=26mm,headheight=15pt]{geometry}
\usepackage[T1]{fontenc}
\usepackage[utf8]{inputenc}
\IfFileExists{libertinus.sty}{\usepackage{libertinus}}{\usepackage{lmodern}}
\usepackage{microtype}
\usepackage{amsmath,amssymb,amsthm,mathtools,mathrsfs,bm}
\usepackage{booktabs,longtable,tabularx,array}
\usepackage{graphicx}
\usepackage{xcolor}
\usepackage{enumitem}
\usepackage{fancyhdr}
\usepackage{xurl}
\usepackage{listings}
\usepackage[most]{tcolorbox}
\usepackage{hyperref}
\usepackage[nameinlink,noabbrev,capitalise]{cleveref}

\definecolor{linkblue}{RGB}{24,72,124}
\definecolor{deepgreen}{RGB}{23,102,69}
\definecolor{deepred}{RGB}{140,42,42}
\definecolor{deepgold}{RGB}{122,91,19}
\definecolor{shadegray}{RGB}{246,247,249}
\definecolor{rulegray}{RGB}{195,201,208}
\definecolor{palered}{RGB}{251,236,236}
\definecolor{palegold}{RGB}{251,245,230}
\definecolor{palegreen}{RGB}{233,245,243}
\definecolor{paleblue}{RGB}{236,243,251}

\hypersetup{
  hypertexnames=false,
  colorlinks=true,
  linkcolor=linkblue,
  citecolor=deepgreen,
  urlcolor=linkblue,
  pdftitle={Asymptotic Inversion of the Subfactorial},
  pdfauthor={Vladimir Reshetnikov},
  pdfsubject={Bell-number tails, inverse-gamma anchoring, and transseries inversion},
  pdfkeywords={subfactorial, derangement, Bell numbers, inverse gamma, Lambert W, transseries, Lagrange-Burmann inversion}
}

\pagestyle{fancy}
\fancyhf{}
\fancyhead[L]{\small Inverse subfactorial transseries}
\fancyhead[R]{\small\nouppercase{\leftmark}}
\fancyfoot[C]{\thepage}
\renewcommand{\headrulewidth}{0.4pt}

\setcounter{tocdepth}{2}
\setcounter{secnumdepth}{3}
\setlist[itemize]{leftmargin=1.5em,itemsep=0.25ex,topsep=0.6ex}
\setlist[enumerate]{leftmargin=1.9em,itemsep=0.35ex,topsep=0.6ex}
\allowdisplaybreaks[2]
\setlength{\emergencystretch}{2em}

\theoremstyle{plain}
\newtheorem{theorem}{Theorem}[section]
\newtheorem{proposition}[theorem]{Proposition}
\newtheorem{lemma}[theorem]{Lemma}
\newtheorem{corollary}[theorem]{Corollary}
\theoremstyle{definition}
\newtheorem{definition}[theorem]{Definition}
\newtheorem{algorithm}[theorem]{Algorithm}
\newtheorem{example}[theorem]{Example}
\theoremstyle{remark}
\newtheorem{remark}[theorem]{Remark}
\newtheorem{convention}[theorem]{Convention}

\crefname{algorithm}{algorithm}{algorithms}
\Crefname{algorithm}{Algorithm}{Algorithms}

\newtcolorbox{mainbox}[1][]{%
  enhanced,breakable,colback=paleblue,colframe=linkblue,
  boxrule=0.7pt,arc=1.2mm,left=2mm,right=2mm,top=1.4mm,bottom=1.4mm,
  title={Main result},fonttitle=\bfseries,#1}
\newtcolorbox{warningbox}[1][]{%
  enhanced,breakable,colback=palered,colframe=deepred,
  boxrule=0.6pt,arc=1.2mm,left=2mm,right=2mm,top=1.2mm,bottom=1.2mm,
  title={Scale warning},fonttitle=\bfseries,#1}
\newtcolorbox{methodbox}[1][]{%
  enhanced,breakable,colback=palegreen,colframe=deepgreen,
  boxrule=0.6pt,arc=1.2mm,left=2mm,right=2mm,top=1.2mm,bottom=1.2mm,
  title={Reusable method},fonttitle=\bfseries,#1}
\newtcolorbox{summarybox}[1][]{%
  enhanced,breakable,colback=palegold,colframe=deepgold,
  boxrule=0.6pt,arc=1.2mm,left=2mm,right=2mm,top=1.2mm,bottom=1.2mm,
  title={Summary},fonttitle=\bfseries,#1}

\lstdefinestyle{pythonstyle}{%
  language=Python,
  basicstyle=\ttfamily\small,
  columns=fullflexible,
  frame=single,
  rulecolor=\color{rulegray},
  backgroundcolor=\color{shadegray},
  showstringspaces=false,
  keepspaces=true,
  breaklines=true,
  xleftmargin=0.4em,xrightmargin=0.4em,
  aboveskip=0.8em,belowskip=0.8em
}
\lstset{style=pythonstyle}

\newcommand{\ee}{\mathrm{e}}
\newcommand{\dd}{\mathop{}\!\mathrm{d}}
\newcommand{\R}{\mathbb{R}}
\newcommand{\C}{\mathbb{C}}
\newcommand{\N}{\mathbb{N}}
\newcommand{\W}{\operatorname{W}}
\newcommand{\Dfrak}{\mathfrak{D}}
\newcommand{\Xfrak}{\mathfrak{X}}
\newcommand{\Tcal}{\mathcal{T}}
\newcommand{\Lcal}{\mathcal{L}}
\newcommand{\Acal}{\mathcal{A}}
\newcommand{\Bell}{\mathrm{B}}
\newcommand{\poch}[2]{\left(#1\right)_{#2}}
\newcommand{\abs}[1]{\left\lvert#1\right\rvert}
\newcommand{\norm}[1]{\left\lVert#1\right\rVert}
\newcommand{\bigO}{\mathcal{O}}
\DeclareMathOperator{\RePart}{Re}
\DeclareMathOperator{\Res}{Res}
\DeclareMathOperator{\sgn}{sgn}
\DeclareMathOperator{\diag}{diag}

\begin{document}

\title{\textbf{Asymptotic Inversion of the Subfactorial}\\[0.35em]
  \Large Bell-Number Tails, Inverse-Gamma Anchoring, and a Calculus for Exponentially Separated Transseries}
\author{Vladimir Reshetnikov\\
  \small ProveIt research note}
\date{3 September 2026}
\maketitle

\begin{abstract}
Let \(D_n={!n}\) denote the number of derangements of \(n\) objects.  The
integer sequence itself has no unique continuous inverse, so the first task is
to specify an interpolation.  We use the real median of the two lateral
incomplete-gamma continuations,
\[
 \Dfrak(x)=\frac{\Gamma(x+1)}{\ee}+\cos(\pi x)J(x),
 \qquad
 J(x)=\int_0^1 \ee^{-t}(1-t)^x\,\dd t,
 \qquad x>-1.
\]
It satisfies \(\Dfrak(n)=D_n\), is strictly increasing for \(x\ge3\), and
therefore has a principal inverse \(\Xfrak(y)\) for \(y\ge2\).  The coefficient
\(J\) has the convergent inverse-factorial representation
\[
 J(x)=\sum_{r\ge0}\frac{(-1)^r}{(x+1)_{r+1}}
\]
and the Bell-number expansion
\[
 J(x)\sim\sum_{k\ge1}(-1)^{k-1}\frac{\Bell_k}{x^k},
 \qquad
 0<(-1)^M\!\left(J(x)-\sum_{k=1}^M(-1)^{k-1}\frac{\Bell_k}{x^k}\right)
 \le \frac{\Bell_{M+1}}{x^{M+1}}.
\]
This extends the known integer formula to every real \(x>0\), with its sharp
sign and constant.

Write \(G(x)=\Gamma(x+1)/\ee\) and let \(g=G^{-1}\) be its increasing large
branch.  With
\[
 A=\log\!\frac{\ee y}{\sqrt{2\pi}},\qquad
 w=\W_0(A/\ee),\qquad X=\frac{A}{w}=\ee^{w+1},\qquad L=\log X,
\]
we derive the outer inverse-gamma expansion
\[
 g(y)=X-\frac12+\sum_{m\ge1}\frac{a_m(L)}{X^{2m-1}},
\]
including an explicit triangular recurrence and the first five rational
coefficient blocks.  The subfactorial correction is beyond every algebraic
order in \(X^{-1}\).  Conjugating by \(G\) and applying the
Lagrange--B\"urmann formula gives the all-orders exact-coefficient transseries
\[
 \boxed{\quad
 \Xfrak(y)=g(y)+\sum_{n\ge1}\frac{(-1)^n}{n!}
 \left(\frac{\dd}{\dd y}\right)^{n-1}
 \bigl[g'(y)q(g(y))^n\bigr],\qquad q(x)=\cos(\pi x)J(x).
 \quad}
\]
Its first sector is
\[
 -\frac{\cos(\pi g(y))}{y\,\psi(g(y)+1)}
 \left(\frac1{g(y)}-\frac2{g(y)^2}+\frac5{g(y)^3}-\cdots\right),
\]
and the \(n\)-th sector has size
\(\bigO(y^{-n}g(y)^{-n}/\log g(y))\) for fixed \(n\).  We also prove a
universal operator formula for inverting
\(G(x)+\sum_j\sigma_jG(x)^{\alpha_j}h_j(x)\), with sector grade
\(\sum_jm_j(1-\alpha_j)\).  The article treats interpolation ambiguity,
nested late-term effects of the Bell expansion, convergence and remainder
conditions, symbolic generation, numerical validation, residual certificates,
and recovery of the integer generalized inverse.
\end{abstract}

\noindent\textbf{Keywords.}
Subfactorial; derangement; Bell numbers; inverse gamma function; Lambert
\(W\); logarithmic transseries; Lagrange--B\"urmann inversion; asymptotic
reversion; exponentially small corrections; incomplete gamma function.

{\small\tableofcontents}
\clearpage

\section{The inverse problem and the answer in three layers}
\label{sec:introduction}

The subfactorial or derangement number is
\begin{equation}
 D_n={!n}=n!\sum_{k=0}^{n}\frac{(-1)^k}{k!},
 \qquad n\in\N_0.
 \label{eq:derangement-def}
\end{equation}
For \(n\ge1\), the familiar estimate \(D_n\approx n!/\ee\) is extraordinarily accurate:
indeed, \(D_n\) is the nearest integer to \(n!/\ee\).  Hassani's integral
identity refines this by isolating the signed difference and leads to the
Bell-number expansion
\begin{equation}
 D_n=\frac{n!}{\ee}
 +\sum_{k=1}^{M}(-1)^{n+k-1}\frac{\Bell_k}{n^k}
 +\bigO(n^{-M-1}),
 \label{eq:hassani-intro}
\end{equation}
with implied constant at most \(\Bell_{M+1}\); see
\cite{Hassani2020}.  The purpose of this article is to invert the full object,
not merely its factorial leading term.

There are three asymptotic layers.
\begin{enumerate}
\item The dominant equation \(\Gamma(x+1)/\ee=y\) is inverted by a
Lambert-\(W\)-anchored logarithmic expansion.  This is the \emph{zero sector}.
\item The Bell tail contributes a correction of order
\(y^{-1}x^{-1}(\log x)^{-1}\), oscillating through \(\cos(\pi x)\).  It is
smaller than every power of the zero-sector variable \(X^{-1}\).
\item Repeated reversion creates sectors of order \(y^{-2},y^{-3},\ldots\).
Each sector is a finite Fourier polynomial whose coefficients have
power-logarithmic expansions generated from the Bell numbers and polygamma
functions.
\end{enumerate}
These layers must not be collapsed into a single ordinary Poincar\'e series.

\begin{warningbox}[title={An integer sequence does not determine a unique continuous inverse}]
Knowing \(D_n\) at integer arguments does not determine a function between the
integers.  If \(R(x)\) is any sufficiently small smooth function, then
\[
 \widetilde{\Dfrak}(x)=\Dfrak(x)+\sin(\pi x)R(x)
\]
has the same values at every integer but a different beyond-all-orders inverse
transseries.  Consequently, the zero-sector inverse-gamma expansion is robust,
whereas the oscillatory sectors depend on the chosen interpolation.  We make
the interpolation explicit before speaking of an inverse.
\end{warningbox}

\begin{mainbox}
Define
\begin{equation}
 \Dfrak(x)=G(x)+q(x),\qquad
 G(x)=\frac{\Gamma(x+1)}{\ee},\qquad
 q(x)=\cos(\pi x)J(x),
 \label{eq:main-interpolation}
\end{equation}
where
\begin{equation}
 J(x)=\int_0^1\ee^{-t}(1-t)^x\,\dd t.
 \label{eq:J-def-intro}
\end{equation}
Then \(\Dfrak(n)=D_n\) for every \(n\ge0\), and \(\Dfrak\) is strictly
increasing on \([3,\infty)\).  Let
\begin{equation}
 \Xfrak=\left(\Dfrak\vert_{[3,\infty)}\right)^{-1}:[2,\infty)\longrightarrow[3,\infty).
 \label{eq:X-def}
\end{equation}
If \(g=G^{-1}\) is the increasing large branch, then
\begin{equation}
 \Xfrak(y)=g(y)+\sum_{n\ge1}\Tcal_n(y),
 \qquad
 \Tcal_n(y)=\frac{(-1)^n}{n!}
 \left(\frac{\dd}{\dd y}\right)^{n-1}
 \left[g'(y)q(g(y))^n\right].
 \label{eq:main-LB}
\end{equation}
For every fixed \(N\), this is an asymptotic expansion through sector \(N\);
under the analytic hypotheses stated in \Cref{sec:validity}, it is the Taylor
series at perturbation parameter zero and converges at parameter one for all
sufficiently large \(y\).
\end{mainbox}

The rest of the paper derives \eqref{eq:main-LB}, expands all its ingredients,
and generalizes it to a broad class of rapidly increasing dominant functions.

\section{A canonical real interpolation of the subfactorial}
\label{sec:interpolation}

\subsection{From Hassani's integral to a median continuation}

Hassani's identity can be written
\begin{equation}
 D_n=\frac{n!}{\ee}+\frac{(-1)^n}{\ee}L_n(\ee),
 \qquad
 L_n(\ee)=\int_1^{\ee}(\log t)^n\,\dd t.
 \label{eq:hassani-integral}
\end{equation}
The substitutions \(t=\ee^u\) and then \(u=1-s\) give
\begin{equation}
 \frac{1}{\ee}L_x(\ee)
 =\frac1\ee\int_0^1u^x\ee^u\,\dd u
 =\int_0^1\ee^{-s}(1-s)^x\,\dd s
 =J(x),
 \label{eq:L-to-J}
\end{equation}
which is real and positive for \(x>-1\).  Replacing \((-1)^n\) by
\(\cos(\pi x)\) gives \eqref{eq:main-interpolation}.

There is also a useful complex interpretation.  Continue the upper incomplete
gamma integral to the point \(-1\) along the upper or lower bank of the
negative real axis.  For \(\Re z>-1\), splitting the contour at the origin
gives the two lateral continuations
\begin{equation}
 \Dfrak_{\pm}(z)=\frac{1}{\ee}\Gamma(z+1,-1)_{\pm}
 =\frac{\Gamma(z+1)}{\ee}+\ee^{\pm i\pi z}J(z).
 \label{eq:lateral-D}
\end{equation}
Their real median on the real axis is
\begin{equation}
 \Dfrak(x)=\frac12\left(\Dfrak_+(x)+\Dfrak_-(x)\right)
 =G(x)+\cos(\pi x)J(x).
 \label{eq:median-D}
\end{equation}
At integer \(n\), both lateral values coincide because
\(\ee^{\pm i\pi n}=(-1)^n\).

\begin{proposition}[Interpolation property]
\label{prop:interpolation}
For every integer \(n\ge0\),
\begin{equation}
 \Dfrak(n)=D_n.
 \label{eq:interpolation-property}
\end{equation}
Moreover, \(\Dfrak\) is real analytic on \((-1,\infty)\).
\end{proposition}
\begin{proof}
Equation \eqref{eq:interpolation-property} follows immediately from
\eqref{eq:hassani-integral}--\eqref{eq:L-to-J}.  Real analyticity of \(J\)
on \(x>-1\) follows by differentiating under the integral sign on compact
subintervals; the gamma and cosine factors are analytic there as well.
\end{proof}

\subsection{Monotonicity and the principal inverse}

The following elementary bounds are enough to isolate an unambiguous large
branch.

\begin{lemma}[Bounds for the Bell tail and its derivatives]
\label{lem:J-derivative-bound}
For \(x>0\) and every integer \(r\ge0\),
\begin{equation}
 \abs{J^{(r)}(x)}\le\frac{r!}{(x+1)^{r+1}}.
 \label{eq:J-derivative-bound}
\end{equation}
In particular,
\begin{equation}
 \abs{q'(x)}\le\frac{\pi}{x+1}+\frac{1}{(x+1)^2}.
 \label{eq:qprime-bound}
\end{equation}
\end{lemma}
\begin{proof}
Set \(s=1-\ee^{-u}\) in \eqref{eq:J-def-intro}.  Then
\begin{equation}
 J(x)=\int_0^\infty\ee^{-xu}K(u)\,\dd u,
 \qquad
 K(u)=\ee^{-u+\ee^{-u}-1}.
 \label{eq:J-Laplace-preview}
\end{equation}
Since \(K(u)\le\ee^{-u}\) for \(u\ge0\), differentiation under the integral
gives
\[
 \abs{J^{(r)}(x)}
 \le\int_0^\infty u^r\ee^{-(x+1)u}\,\dd u
 =\frac{r!}{(x+1)^{r+1}}.
\]
Finally,
\(q'=-\pi\sin(\pi x)J+\cos(\pi x)J'\), which yields
\eqref{eq:qprime-bound}.
\end{proof}

\begin{theorem}[Existence of the principal inverse]
\label{thm:monotonicity}
The function \(\Dfrak\) is strictly increasing on \([3,\infty)\), maps that
interval onto \([2,\infty)\), and therefore has the inverse \(\Xfrak\) defined
in \eqref{eq:X-def}.
\end{theorem}
\begin{proof}
The dominant derivative is
\begin{equation}
 G'(x)=G(x)\psi(x+1),
 \label{eq:Gprime-digamma}
\end{equation}
where \(\psi\) is the digamma function.  For \(x\ge3\), both
\(\Gamma(x+1)\) and \(\psi(x+1)\) are increasing, and hence
\[
 G'(x)\ge\frac{6}{\ee}\psi(4)
 =\frac{6}{\ee}\left(\frac{11}{6}-\gamma\right)>2.7.
\]
By \eqref{eq:qprime-bound},
\[
 \abs{q'(x)}\le\frac\pi4+\frac1{16}<0.85.
\]
Thus \(\Dfrak'(x)=G'(x)+q'(x)>0\) on \([3,\infty)\).  Since
\(\Dfrak(3)=D_3=2\) and \(\Dfrak(x)\to\infty\), the image is
\([2,\infty)\).
\end{proof}

\begin{corollary}[Integer generalized inverse]
\label{cor:integer-inverse}
For \(y\ge2\), define
\begin{equation}
 N(y)=\max\{n\in\N:n\ge3,\ D_n\le y\}.
 \label{eq:integer-generalized-inverse}
\end{equation}
Then
\begin{equation}
 \boxed{\quad N(y)=\lfloor\Xfrak(y)\rfloor.\quad}
 \label{eq:floor-inverse}
\end{equation}
In particular, \(\Xfrak(D_n)=n\) for every \(n\ge3\).
\end{corollary}
\begin{proof}
Strict monotonicity gives
\(D_n=\Dfrak(n)\le y<D_{n+1}=\Dfrak(n+1)\) exactly when
\(n\le\Xfrak(y)<n+1\).
\end{proof}

\section{The Bell-number coefficient: exact forms, asymptotics, and late terms}
\label{sec:Bell-tail}

The coefficient \(J(x)\) is the bridge between the discrete derangement
identity and the continuous inverse transseries.  It is advantageous to retain
several equivalent representations because they serve different purposes:
the integral gives positivity and derivative bounds, the inverse-factorial
series is convergent and numerically stable, and the Bell series exposes the
power asymptotics.

\subsection{A convergent inverse-factorial series}

\begin{theorem}[Exact representations of \(J\)]
\label{thm:J-exact}
For every real \(x>-1\),
\begin{align}
 J(x)
 &=\int_0^1\ee^{-t}(1-t)^x\,\dd t,
 \label{eq:J-int}\\
 &=\frac1\ee\int_0^1\ee^u u^x\,\dd u,
 \label{eq:J-u-int}\\
 &=\sum_{r=0}^{\infty}\frac{(-1)^r}{(x+1)_{r+1}},
 \label{eq:J-inverse-factorial}\\
 &=\int_0^\infty\ee^{-xu}\ee^{-u+\ee^{-u}-1}\,\dd u.
 \label{eq:J-Laplace}
\end{align}
The series \eqref{eq:J-inverse-factorial} converges absolutely.  For \(x>0\),
its alternating truncation error is bounded by the first omitted term:
\begin{equation}
 \abs{J(x)-\sum_{r=0}^{R-1}\frac{(-1)^r}{(x+1)_{r+1}}}
 \le \frac{1}{(x+1)_{R+1}}.
 \label{eq:J-inverse-factorial-error}
\end{equation}
\end{theorem}
\begin{proof}
Equivalence of \eqref{eq:J-int} and \eqref{eq:J-u-int} follows from
\(u=1-t\).  Expanding \(\ee^{-t}\) in \eqref{eq:J-int} and integrating beta
integrals yields
\[
 \int_0^1\frac{(-t)^r}{r!}(1-t)^x\,\dd t
 =\frac{(-1)^r}{r!}\,\frac{\Gamma(r+1)\Gamma(x+1)}
 {\Gamma(x+r+2)}
 =\frac{(-1)^r}{(x+1)_{r+1}}.
\]
Absolute convergence follows from comparison with a factorial tail.  For
\(x>0\), the magnitudes decrease strictly, so the alternating-series bound
gives \eqref{eq:J-inverse-factorial-error}.  Finally, the substitution
\(t=1-\ee^{-u}\) gives \eqref{eq:J-Laplace}.
\end{proof}

The recurrence
\begin{equation}
 t_0=\frac1{x+1},\qquad
 t_{r+1}=-\frac{t_r}{x+r+2},\qquad
 J(x)=\sum_{r\ge0}t_r,
 \label{eq:J-term-recurrence}
\end{equation}
computes \(J\) without large intermediate quantities.  It is preferable to
the divergent Bell expansion when many digits are needed.

\subsection{Why Bell numbers appear}

Set
\begin{equation}
 K(u)=\ee^{-u+\ee^{-u}-1}.
 \label{eq:K-def}
\end{equation}
The exponential mixture
\begin{equation}
 K(u)=\frac1\ee\sum_{r=0}^{\infty}\frac{\ee^{-(r+1)u}}{r!}
 \label{eq:K-mixture}
\end{equation}
shows that \(K\) is completely monotone on \([0,\infty)\).

\begin{lemma}[Bell derivatives]
\label{lem:K-Bell}
For every integer \(m\ge0\),
\begin{equation}
 K^{(m)}(0)=(-1)^m\Bell_{m+1}.
 \label{eq:K-Bell-derivatives}
\end{equation}
\end{lemma}
\begin{proof}
Differentiate \eqref{eq:K-mixture} termwise:
\[
 K^{(m)}(0)=\frac{(-1)^m}{\ee}
 \sum_{r=0}^{\infty}\frac{(r+1)^m}{r!}
 =\frac{(-1)^m}{\ee}
 \sum_{k=1}^{\infty}\frac{k^{m+1}}{k!}.
\]
Dobinski's formula \(\Bell_j=\ee^{-1}\sum_{k\ge0}k^j/k!\) gives the result.
\end{proof}

\begin{theorem}[Bell expansion with a signed remainder]
\label{thm:Bell-expansion}
For every integer \(M\ge1\) and every real \(x>0\),
\begin{equation}
 J(x)=\sum_{k=1}^{M}(-1)^{k-1}\frac{\Bell_k}{x^k}+R_M(x),
 \label{eq:J-Bell-truncated}
\end{equation}
where
\begin{equation}
 0<(-1)^M R_M(x)\le\frac{\Bell_{M+1}}{x^{M+1}}.
 \label{eq:J-Bell-remainder}
\end{equation}
Consequently,
\begin{equation}
 J(x)\sim\frac1x-\frac2{x^2}+\frac5{x^3}-\frac{15}{x^4}
 +\frac{52}{x^5}-\frac{203}{x^6}+\frac{877}{x^7}-\cdots.
 \label{eq:J-Bell-explicit}
\end{equation}
\end{theorem}
\begin{proof}
Taylor's formula at the origin gives
\[
 K(u)=\sum_{m=0}^{M-1}\frac{K^{(m)}(0)}{m!}u^m+\rho_M(u).
\]
Complete monotonicity implies
\[
 0<(-1)^M\rho_M(u)
 \le \frac{(-1)^M K^{(M)}(0)}{M!}u^M
 =\frac{\Bell_{M+1}}{M!}u^M.
\]
Insert this formula into the Laplace representation \eqref{eq:J-Laplace}.
Because
\(\int_0^\infty\ee^{-xu}u^m\,\dd u=m!/x^{m+1}\),
\Cref{lem:K-Bell} gives \eqref{eq:J-Bell-truncated} and
\eqref{eq:J-Bell-remainder}.
\end{proof}

\begin{corollary}[Hassani's expansion, with the same bound]
\label{cor:Hassani}
For integers \(n\ge1\),
\begin{equation}
 D_n=\frac{n!}{\ee}
 +\sum_{k=1}^{M}(-1)^{n+k-1}\frac{\Bell_k}{n^k}
 +(-1)^nR_M(n),
 \label{eq:Hassani-recovered}
\end{equation}
where \(\abs{R_M(n)}\le\Bell_{M+1}/n^{M+1}\).
\end{corollary}
\begin{proof}
Use \(D_n=n!/\ee+(-1)^nJ(n)\) and
\Cref{thm:Bell-expansion}.
\end{proof}

The derivative expansions follow by termwise differentiation at any fixed
order.

\begin{corollary}[Derivatives of the Bell tail]
\label{cor:J-derivative-expansion}
For each fixed \(r\ge0\),
\begin{equation}
 J^{(r)}(x)\sim
 (-1)^r\sum_{k\ge1}(-1)^{k-1}
 \frac{(k)_r\Bell_k}{x^{k+r}},
 \qquad x\to+\infty,
 \label{eq:J-derivative-expansion}
\end{equation}
where \((k)_r=k(k+1)\cdots(k+r-1)\) is the rising factorial.
\end{corollary}

\subsection{Late terms and the nested transseries warning}
\label{subsec:Bell-late-terms}

The Bell series is asymptotic, not convergent.  This matters when it is placed
inside the exponentially small inverse sectors.  The standard Bell-number
estimate may be written in the form
\begin{equation}
 \Bell_m\sim
 \frac{N^m\ee^{N-m-1}}{\sqrt{1+\log N}},
 \qquad N\log N=m,
 \label{eq:Bell-DLMF-asymptotic}
\end{equation}
with higher corrections available; see \cite{DLMF26,MoserWyman1955}.  The
ratio heuristic places the least term of \(\Bell_m/x^m\) near
\begin{equation}
 m_\star=x\log x,
 \qquad N=x.
 \label{eq:Bell-optimal-index}
\end{equation}
At this index, \eqref{eq:Bell-DLMF-asymptotic} gives
\begin{equation}
 \frac{\Bell_{m_\star}}{x^{m_\star}}
 \asymp \frac{\ee^{-x\log x+x-1}}{\sqrt{1+\log x}}
 \asymp
 \frac{\sqrt{x/\log x}}{G(x)}.
 \label{eq:Bell-least-term-scale}
\end{equation}
More precisely, the quotient of the two displayed leading scales tends to
\(\sqrt{2\pi}/\ee^2\).

\begin{warningbox}[title={The Bell coefficient expansion invades the next exponential sector}]
At \(x=g(y)\), one has \(G(x)=y\).  If the first inverse sector
\(-q/(y\psi)\) is evaluated by optimally truncating the Bell series for \(J\),
then the induced least-term uncertainty is of order
\begin{equation}
 \frac{1}{y\log x}\,
rac{\sqrt{x/\log x}}{y}
 =\bigO\!\left(\frac{\sqrt{x}}{y^2(\log x)^{3/2}}\right).
 \label{eq:Bell-sector-mixing}
\end{equation}
It therefore lies in the \(y^{-2}\) exponential grade and can be larger than
the explicit reversion-generated \(y^{-2}\) block.  A fully expanded
transseries is nested: late terms of one coefficient sector mix with later
exponential sectors.  To preserve clean grading, retain \(J\) exactly through
\eqref{eq:J-int} or the convergent series \eqref{eq:J-inverse-factorial}, and
expand it in Bell numbers only to the accuracy actually required.
\end{warningbox}

This phenomenon is not a defect.  It is the expected interaction between a
power asymptotic expansion and a second, exponentially smaller scale.  The
general calculus in \Cref{sec:general-theory} is therefore formulated first
with exact coefficient functions and only then with their internal
transseries.

\section{The zero sector: inversion of \texorpdfstring{\(\Gamma(x+1)/\ee\)}{Gamma(x+1)/e}}
\label{sec:inverse-gamma}

Let
\begin{equation}
 G(x)=\frac{\Gamma(x+1)}{\ee},
 \qquad g(y)=G^{-1}(y)
 \label{eq:G-g-def}
\end{equation}
where \(g\) denotes the increasing large branch.  We first construct a useful
Lambert-\(W\)-anchored expansion for \(g\).

\subsection{Centered Stirling expansion}

Set
\begin{equation}
 u=x+\frac12.
 \label{eq:u-centered}
\end{equation}
The centered Stirling expansion is
\begin{equation}
 \log\Gamma(x+1)
 \sim u(\log u-1)+\frac12\log(2\pi)
 +\sum_{m\ge1}\frac{\kappa_m}{u^{2m-1}},
 \label{eq:centered-Stirling}
\end{equation}
where
\begin{equation}
 \kappa_m=-\frac{(1-2^{1-2m})\Bell^{\mathrm{Ber}}_{2m}}
 {2m(2m-1)}.
 \label{eq:kappa-def}
\end{equation}
Here \(\Bell^{\mathrm{Ber}}_{2m}\) denotes a Bernoulli number; the superscript
is included only to prevent confusion with the Bell numbers \(\Bell_k\).
The first terms are
\begin{equation}
 \kappa_1=-\frac1{24},\qquad
 \kappa_2=\frac7{2880},\qquad
 \kappa_3=-\frac{31}{40320},\qquad
 \kappa_4=\frac{127}{215040}.
 \label{eq:kappa-first}
\end{equation}
This centered form, obtained from the standard shifted gamma expansion, has
only odd inverse powers and is especially convenient for reversion; see
\cite{DLMF5,Olver1997}.

The equation \(G(x)=y\) becomes
\begin{equation}
 u(\log u-1)+\sum_{m\ge1}\frac{\kappa_m}{u^{2m-1}}
 =A,
 \qquad
 A=\log\!\frac{\ee y}{\sqrt{2\pi}}.
 \label{eq:inverse-gamma-equation}
\end{equation}
Ignore the inverse powers temporarily and solve
\begin{equation}
 X(\log X-1)=A.
 \label{eq:X-leading-equation}
\end{equation}
With the principal Lambert function,
\begin{equation}
 w=\W_0(A/\ee),\qquad
 X=\frac{A}{w}=\ee^{w+1},\qquad
 L=\log X=w+1.
 \label{eq:AXwL}
\end{equation}
Indeed, \(X(L-1)=Xw=A\).

\subsection{Triangular reversion around the Lambert anchor}

Write
\begin{equation}
 u=X(1+\eta),\qquad
 t=X^{-1},\qquad
 \eta(t)=\sum_{m\ge1}a_m(L)t^{2m}.
 \label{eq:eta-ansatz}
\end{equation}
Substitution in \eqref{eq:inverse-gamma-equation} yields the single formal
identity
\begin{equation}
 0=t^{-1}\left\{\eta(L-1)+(1+\eta)\log(1+\eta)\right\}
 +\sum_{r\ge1}\kappa_rt^{2r-1}(1+\eta)^{-(2r-1)}.
 \label{eq:eta-master}
\end{equation}
At order \(t^{2m-1}\), the unknown \(a_m\) occurs only in the linear term
\(La_m\).  This gives a triangular algorithm.

\begin{algorithm}[Outer inverse-gamma coefficients]
\label{alg:outer-coefficients}
Suppose \(a_1,\ldots,a_{m-1}\) are known.  Put
\[
 \eta_{<m}(t)=\sum_{j=1}^{m-1}a_j(L)t^{2j}
\]
and let \(\rho_m(L)\) be the coefficient of \(t^{2m-1}\) in the right-hand
side of \eqref{eq:eta-master} after replacing \(\eta\) by
\(\eta_{<m}+a_m t^{2m}\), with the term \(La_m\) omitted.  Then
\begin{equation}
 a_m(L)=-\frac{\rho_m(L)}{L}.
 \label{eq:am-recurrence}
\end{equation}
Equivalently, expand \eqref{eq:eta-master} through order \(t^{2m-1}\) and solve
the resulting linear equation for \(a_m\).
\end{algorithm}

\begin{theorem}[Lambert-anchored inverse-gamma expansion]
\label{thm:inverse-gamma-expansion}
As \(y\to+\infty\), with \(A,w,X,L\) defined by \eqref{eq:AXwL},
\begin{equation}
 \boxed{\quad
 g(y)\sim X-\frac12+\sum_{m\ge1}\frac{a_m(L)}{X^{2m-1}}.
 \quad}
 \label{eq:g-outer-main}
\end{equation}
The first five coefficient blocks are
\begin{align}
 a_1(L)&=\frac{1}{24L},
 \label{eq:a1}\\
 a_2(L)&=-\frac{14L^2+10L+5}{5760L^3},
 \label{eq:a2}\\
 a_3(L)&=\frac{2232L^4+1176L^3+714L^2+350L+105}
 {2903040L^5},
 \label{eq:a3}\\
 a_4(L)&=-\frac{1}{1393459200L^7}
 \bigl(822960L^6+292536L^5+136956L^4
 \notag\\[-0.2em]
 &\hspace{8em}{}+68040L^3+32970L^2+12250L+2625\bigr),
 \label{eq:a4}\\
 a_5(L)&=\frac{1}{367873228800L^9}
 \bigl(309052800L^8+77920128L^7+26024592L^6
 \notag\\[-0.2em]
 &\hspace{4.4em}{}+10019592L^5+4516028L^4+2023560L^3
 \notag\\[-0.2em]
 &\hspace{4.4em}{}+812350L^2+242550L+40425\bigr).
 \label{eq:a5}
\end{align}
After truncation through \(a_M/X^{2M-1}\), the asymptotic error is
\begin{equation}
 \bigO\!\left(\frac{1}{L X^{2M+1}}\right).
 \label{eq:g-outer-error}
\end{equation}
Moreover, \(a_m(L)\) is a rational polynomial in \(L^{-1}\), beginning at
order \(L^{-1}\) and ending at order \(L^{-(2m-1)}\).
\end{theorem}
\begin{proof}
Substitution of \eqref{eq:eta-ansatz} into the centered Stirling equation gives
\eqref{eq:eta-master}.  Because
\[
 (1+\eta)\log(1+\eta)=\eta+\frac12\eta^2-\frac16\eta^3+\cdots,
\]
the coefficient multiplying the new \(a_m\) is
\((L-1)+1=L\); all other contributions depend only on previous coefficients.
This proves triangularity and \eqref{eq:am-recurrence}.  Straight expansion
gives \eqref{eq:a1}--\eqref{eq:a5}.  The next omitted term in the centered
Stirling series is \(\bigO(X^{-(2M+1)})\), while the derivative of
\(u(\log u-1)\) at \(u=X\) is \(\log X=L\).  Division by this derivative gives
\eqref{eq:g-outer-error}.  The stated rational structure follows inductively
from \eqref{eq:eta-master}.
\end{proof}

\subsection{Conversion to a pure logarithmic expansion}

The Lambert function itself has the standard expansion
\begin{align}
 \W_0(z)\sim{}&L_1-L_2+\frac{L_2}{L_1}
 +\frac{L_2(-2+L_2)}{2L_1^2}
 \notag\\
 &+\frac{L_2(6-9L_2+2L_2^2)}{6L_1^3}
 +\frac{L_2(-12+36L_2-22L_2^2+3L_2^3)}{12L_1^4}
 +\cdots,
 \label{eq:W-asymptotic}
\end{align}
where \(L_1=\log z\) and \(L_2=\log L_1\); see
\cite{Corless1996,DLMF4}.  Substituting \(z=A/\ee\) into
\eqref{eq:W-asymptotic}, then using \(X=A/W_0(A/\ee)\), turns
\eqref{eq:g-outer-main} into a conventional logarithmic transseries in
\(A,\log A,\log\log A\).  The Lambert-anchored form is usually superior: it
resums the entire leading logarithmic inversion exactly and leaves only odd
powers of \(X^{-1}\).

\begin{remark}[What the zero sector can and cannot see]
\label{rem:zero-sector-blindness}
Since
\[
 y=G(g(y))\asymp \exp(g(y)\log g(y)-g(y)),
\]
one has \(y^{-1}=o(X^{-N})\) for every fixed \(N\).  Thus no number of terms in
\eqref{eq:g-outer-main} can reveal the Bell correction.  The latter is a
separate exponential sector, not a further coefficient in the same
\(X^{-1}\) series.
\end{remark}

\section{The full inverse subfactorial transseries}
\label{sec:full-transseries}

We now restore the additive correction
\(q(x)=\cos(\pi x)J(x)\).  The cleanest derivation is to conjugate the
subfactorial interpolation by the dominant map \(G\).

\subsection{Conjugation to a near-identity problem}

Let
\begin{equation}
 z=G(x),\qquad x=g(z),\qquad Q(z)=q(g(z)).
 \label{eq:conjugation-def}
\end{equation}
Then
\begin{equation}
 \Dfrak(g(z))=z+Q(z).
 \label{eq:near-identity-z}
\end{equation}
If \(h\) denotes the inverse of \(z\mapsto z+Q(z)\), then
\begin{equation}
 \Xfrak(y)=g(h(y)).
 \label{eq:X-conjugated}
\end{equation}
The generalized Lagrange--B\"urmann formula applied to the composition
\(g\circ h\) gives the answer in one line.

\begin{theorem}[Exact-coefficient inverse transseries]
\label{thm:exact-coefficient-transseries}
As a formal series in a perturbation parameter \(\varepsilon\), the inverse of
\begin{equation}
 \Dfrak_\varepsilon(x)=G(x)+\varepsilon q(x)
 \label{eq:D-epsilon}
\end{equation}
has the expansion
\begin{equation}
 \Xfrak_\varepsilon(y)=g(y)+\sum_{n\ge1}\varepsilon^n\Tcal_n(y),
 \label{eq:X-epsilon-series}
\end{equation}
where
\begin{equation}
 \boxed{\quad
 \Tcal_n(y)=\frac{(-1)^n}{n!}
 \left(\frac{\dd}{\dd y}\right)^{n-1}
 \left[g'(y)q(g(y))^n\right].
 \quad}
 \label{eq:Tn-y-form}
\end{equation}
Equivalently, if \(A_1(x)=G'(x)\), then
\begin{equation}
 \boxed{\quad
 \Tcal_n(y)=\frac{(-1)^n}{n!}
 \left[
 \left(A_1(x)^{-1}\frac{\dd}{\dd x}\right)^{n-1}
 \left(\frac{q(x)^n}{A_1(x)}\right)
 \right]_{x=g(y)}.
 \quad}
 \label{eq:Tn-x-form}
\end{equation}
Setting \(\varepsilon=1\) gives the inverse subfactorial transseries
\eqref{eq:main-LB}.
\end{theorem}
\begin{proof}
In the conjugated variable, the perturbed equation is
\[
 y=z+\varepsilon Q(z),\qquad Q(z)=q(g(z)).
\]
The Lagrange--B\"urmann composition formula gives
\[
 g(h_\varepsilon(y))
 =g(y)+\sum_{n\ge1}\frac{(-\varepsilon)^n}{n!}
 \left(\frac{\dd}{\dd y}\right)^{n-1}
 \left[g'(y)Q(y)^n\right],
\]
which is \eqref{eq:Tn-y-form}.  Since
\(g'(y)=1/G'(g(y))\) and
\(\dd/\dd y=G'(x)^{-1}\dd/\dd x\) at \(x=g(y)\), one obtains
\eqref{eq:Tn-x-form}.
\end{proof}

Formula \eqref{eq:Tn-x-form} is an algorithm, not merely an existence
statement.  It differentiates only known functions and automatically produces
the correct nonlinear combinations of \(q\), gamma derivatives, and
polygamma functions.

\subsection{The first three exponential sectors}

For compactness, write
\begin{equation}
 A_1=G',\qquad A_2=G'',\qquad A_3=G'''.
 \label{eq:A123}
\end{equation}
Direct expansion of \eqref{eq:Tn-x-form} gives
\begin{align}
 \Tcal_1={}&-\frac{q}{A_1},
 \label{eq:T1-general}\\
 \Tcal_2={}&\frac{qq'}{A_1^2}-\frac{q^2A_2}{2A_1^3},
 \label{eq:T2-general}\\
 \Tcal_3={}&-\frac{q^2q''}{2A_1^3}
 -\frac{q(q')^2}{A_1^3}
 +\frac{q^3A_3}{6A_1^4}
 +\frac{3q^2q'A_2}{2A_1^4}
 -\frac{q^3A_2^2}{2A_1^5},
 \label{eq:T3-general}
\end{align}
where every quantity on the right is evaluated at \(x=g(y)\).

Introduce the polygamma abbreviations
\begin{equation}
 p=\psi(x+1),\qquad
 p_1=\psi^{(1)}(x+1),\qquad
 p_2=\psi^{(2)}(x+1).
 \label{eq:p-polygamma}
\end{equation}
At \(x=g(y)\), \(G(x)=y\), and therefore
\begin{align}
 G'&=yp,
 \label{eq:G1-poly}\\
 G''&=y(p^2+p_1),
 \label{eq:G2-poly}\\
 G'''&=y(p^3+3pp_1+p_2).
 \label{eq:G3-poly}
\end{align}
Consequently,
\begin{equation}
 \boxed{\quad \Tcal_1(y)=-\frac{q(x)}{y\,p},\qquad x=g(y).\quad}
 \label{eq:T1-polygamma}
\end{equation}
The next two blocks are
\begin{align}
 \Tcal_2(y)
 ={}&\frac1{y^2}\left[
 \frac{qq'}{p^2}
 -\frac{q^2(p^2+p_1)}{2p^3}
 \right]_{x=g(y)},
 \label{eq:T2-polygamma}\\
 \Tcal_3(y)
 ={}&\frac1{y^3}\left[
 -\frac{q^2q''}{2p^3}-\frac{q(q')^2}{p^3}
 +\frac{q^3(p^3+3pp_1+p_2)}{6p^4}
 \right.\notag\\[-0.2em]
 &\left.\hspace{3.5em}
 +\frac{3q^2q'(p^2+p_1)}{2p^4}
 -\frac{q^3(p^2+p_1)^2}{2p^5}
 \right]_{x=g(y)}.
 \label{eq:T3-polygamma}
\end{align}
For actual computation one may use
\begin{align}
 q&=cJ,
 \label{eq:q-cJ}\\
 q'&=-\pi sJ+cJ',
 \label{eq:qp-cJ}\\
 q''&=-\pi^2cJ-2\pi sJ'+cJ'',
 \label{eq:qpp-cJ}
\end{align}
where \(c=\cos(\pi x)\), \(s=\sin(\pi x)\).

\subsection{The explicit first Bell--Fourier sector}

The digamma expansion is
\begin{equation}
 \psi(x+1)\sim
 \log x+\frac1{2x}-\frac1{12x^2}+\frac1{120x^4}
 -\frac1{252x^6}+\cdots.
 \label{eq:digamma-expansion}
\end{equation}
Combining it with \eqref{eq:J-Bell-explicit} gives a power-logarithmic
coefficient expansion for the first exponential sector.  Let
\(\ell=\log x\).  Then
\begin{align}
 \frac{J(x)}{\psi(x+1)}\sim{}&
 \frac{1}{x\ell}
 -\frac{4\ell+1}{2x^2\ell^2}
 +\frac{60\ell^2+13\ell+3}{12x^3\ell^3}
 \notag\\
 &-\frac{360\ell^3+64\ell^2+14\ell+3}{24x^4\ell^4}
 \notag\\
 &+\frac{37440\ell^4+5694\ell^3+1025\ell^2+225\ell+45}
 {720x^5\ell^5}
 +\bigO\!\left(\frac1{x^6\ell}\right).
 \label{eq:J-over-psi}
\end{align}
Thus, with \(x=g(y)\),
\begin{align}
 \Xfrak(y)={}&g(y)
 -\frac{\cos(\pi g(y))}{y}
 \left[
 \frac{1}{g\log g}
 -\frac{4\log g+1}{2g^2(\log g)^2}
 \right.\notag\\[-0.2em]
 &\left.\hspace{8em}
 +\frac{60(\log g)^2+13\log g+3}{12g^3(\log g)^3}
 -\cdots
 \right]
 +\Tcal_2(y)+\Tcal_3(y)+\cdots.
 \label{eq:X-first-explicit}
\end{align}
Here and only here, inside the square brackets, \(g\) abbreviates \(g(y)\).
Replacing \(g(y)\) by the Lambert-anchored series
\eqref{eq:g-outer-main} produces a transseries entirely in
\(W_0(A/\ee)\), inverse powers of \(X\), powers of \(y^{-1}\), and the Fourier
phase \(\ee^{\pm i\pi g(y)}\).

The leading form of the first three sectors is especially simple:
\begin{equation}
 \Tcal_n(y)\sim-\frac{q(g(y))^n}{n\,y^n\psi(g(y)+1)}
 \sim-\frac{\cos^n(\pi g(y))}
 {n\,y^n g(y)^n\log g(y)},
 \label{eq:Tn-leading}
\end{equation}
for each fixed \(n\), away from accidental zeros of the displayed leading
coefficient.  The second equivalence uses \(J(x)\sim x^{-1}\).

\subsection{Sector size and Fourier support}

\begin{theorem}[Subfactorial sector structure]
\label{thm:sector-structure}
For each fixed \(n\ge1\),
\begin{equation}
 \Tcal_n(y)=\bigO_n\!\left(
 \frac{1}{y^n g(y)^n\log g(y)}
 \right),
 \qquad y\to+\infty.
 \label{eq:Tn-size}
\end{equation}
After expanding \(J\), the polygamma functions, and the zero sector, the
coefficient of \(y^{-n}\) is a finite Fourier polynomial in
\(\ee^{i\pi g(y)}\) with frequencies
\begin{equation}
 -n,-n+2,\ldots,n-2,n.
 \label{eq:frequency-support}
\end{equation}
Each Fourier coefficient is a power-logarithmic series in \(g(y)^{-1}\) whose
coefficients lie in the ring generated over \(\mathbb{Q}\) by Bell numbers,
Bernoulli numbers, and powers of \(\pi\).
\end{theorem}
\begin{proof}
By \Cref{lem:J-derivative-bound}, every fixed derivative of \(q\) is
\(\bigO_n(x^{-1})\).  Moreover,
\(G'=G\psi\) and \(\psi(x+1)\sim\log x\).  Formula
\eqref{eq:Tn-x-form} therefore contributes one factor \(G^{-1}=y^{-1}\) per
application and starts with \(q^n=\bigO(x^{-n})\), yielding
\eqref{eq:Tn-size}.  Writing
\[
 q(x)=\frac12J(x)\left(\ee^{i\pi x}+\ee^{-i\pi x}\right)
\]
shows that \(q^n\) has precisely the frequencies in
\eqref{eq:frequency-support}.  Differentiation changes coefficients but not
frequencies.  Finally, the Bell expansion of \(J\), the Bernoulli expansion of
the polygamma functions, and differentiation generate the stated coefficient
ring.
\end{proof}

\begin{remark}[Lateral rather than median interpolation]
For \(\Dfrak_\pm=G+\ee^{\pm i\pi x}J\), the \(n\)-th sector carries a single
frequency \(\pm n\pi\).  The real median combines the two lateral modes, and
multinomial products produce the finite frequency fan
\eqref{eq:frequency-support}.
\end{remark}

\section{A general inversion calculus for rapidly increasing transseries}
\label{sec:general-theory}

The preceding calculation is an instance of a universal pattern.  A rapidly
increasing dominant map is inverted first; additive or relatively smaller
transseries are then reverted by a near-identity Lagrange expansion.  The
logarithmic derivative of the dominant map controls every coefficient block.

\subsection{The one-perturbation operator formula}

Let \(G\) be a strictly increasing function on a large real interval and set
\begin{equation}
 \lambda(x)=\frac{G'(x)}{G(x)}=\frac{\dd}{\dd x}\log G(x).
 \label{eq:lambda-def}
\end{equation}
The elementary identity
\begin{equation}
 \frac1{G\lambda}\frac{\dd}{\dd x}\left(G^{-k}f\right)
 =G^{-(k+1)}\left(\lambda^{-1}\frac{\dd}{\dd x}-k\right)f
 \label{eq:operator-factorization}
\end{equation}
is the source of the block structure.  Define
\begin{equation}
 \Lcal_k=\lambda^{-1}\frac{\dd}{\dd x}-k.
 \label{eq:Lk-def}
\end{equation}

\begin{theorem}[Universal additive-block formula]
\label{thm:universal-additive}
Let
\begin{equation}
 F(x)=G(x)+H(x),
 \label{eq:F-GH}
\end{equation}
and let \(g=G^{-1}\).  Formally, and asymptotically under the hypotheses of
\Cref{sec:validity},
\begin{equation}
 F^{-1}(y)=g(y)+\sum_{n\ge1}\frac{U_n(g(y))}{y^n},
 \label{eq:universal-additive-series}
\end{equation}
where
\begin{equation}
 \boxed{\quad
 U_n(x)=\frac{(-1)^n}{n!}
 \Lcal_{n-1}\Lcal_{n-2}\cdots\Lcal_1
 \left(\frac{H(x)^n}{\lambda(x)}\right).
 \quad}
 \label{eq:Un-formula}
\end{equation}
For \(n=1\), the operator product is empty.  The first blocks are
\begin{align}
 U_1&=-\frac{H}{\lambda},
 \label{eq:U1}\\
 U_2&=\frac12\left(\lambda^{-1}D_x-1\right)
 \left(\frac{H^2}{\lambda}\right),
 \label{eq:U2}\\
 U_3&=-\frac16\left(\lambda^{-1}D_x-2\right)
 \left(\lambda^{-1}D_x-1\right)
 \left(\frac{H^3}{\lambda}\right).
 \label{eq:U3}
\end{align}
\end{theorem}
\begin{proof}
Apply \eqref{eq:Tn-x-form} with \(q=H\) and
\(A_1=G'=G\lambda\).  The expression inside the iterated derivative is
\[
 \frac{H^n}{G\lambda}=G^{-1}\frac{H^n}{\lambda}.
\]
Repeated use of \eqref{eq:operator-factorization}, with
\(k=1,2,\ldots,n-1\), produces
\[
 G^{-n}\Lcal_{n-1}\cdots\Lcal_1(H^n/\lambda).
\]
At \(x=g(y)\), \(G(x)=y\), which proves
\eqref{eq:universal-additive-series}--\eqref{eq:Un-formula}.
\end{proof}

For the subfactorial, \(H=q\) and
\begin{equation}
 \lambda(x)=\psi(x+1).
 \label{eq:lambda-subfactorial}
\end{equation}
Thus \eqref{eq:Un-formula} is exactly the compact form of the full
subfactorial sector calculus.

\subsection{Several relative scales and a sector semigroup}

The same factorization handles perturbations that carry powers of the
dominant function.  Consider
\begin{equation}
 F(x)=G(x)+\sum_{j=1}^{s}\sigma_jG(x)^{\alpha_j}h_j(x),
 \qquad \alpha_j<1.
 \label{eq:multi-perturbation}
\end{equation}
Let \(\bm m=(m_1,\ldots,m_s)\in\N_0^s\setminus\{0\}\), and define
\begin{equation}
 n=\abs{\bm m}=\sum_jm_j,
 \qquad
 a=\bm m\!\cdot\!\bm\alpha=\sum_jm_j\alpha_j,
 \qquad
 h_{\bm m}=\prod_jh_j^{m_j},
 \qquad
 \bm m!=\prod_jm_j!.
 \label{eq:multiindex-defs}
\end{equation}

\begin{theorem}[Multiscale inversion and sector grading]
\label{thm:multiscale}
The formal inverse of \eqref{eq:multi-perturbation} has the expansion
\begin{equation}
 F^{-1}(y)=g(y)+
 \sum_{\bm m\ne0}\bm\sigma^{\bm m}
 y^{-\nu(\bm m)}U_{\bm m}(g(y)),
 \label{eq:multiscale-inverse}
\end{equation}
with sector grade
\begin{equation}
 \boxed{\quad
 \nu(\bm m)=n-a=\sum_{j=1}^{s}m_j(1-\alpha_j)>0.
 \quad}
 \label{eq:sector-grade}
\end{equation}
To compute the coefficient, set
\begin{equation}
 V_0(x)=\frac{h_{\bm m}(x)}{\lambda(x)}
 \label{eq:V0}
\end{equation}
and recursively define
\begin{equation}
 V_{r+1}(x)=
 \left(\lambda(x)^{-1}D_x-(1-a+r)\right)V_r(x),
 \qquad 0\le r\le n-2.
 \label{eq:Vr-recurrence}
\end{equation}
Then
\begin{equation}
 \boxed{\quad
 U_{\bm m}(x)=\frac{(-1)^n}{\bm m!}V_{n-1}(x).
 \quad}
 \label{eq:Umulti}
\end{equation}
For \(n=1\), \(V_{n-1}=V_0\).
\end{theorem}
\begin{proof}
In the \(n\)-th Lagrange block, the multinomial expansion of the perturbation
gives
\[
 \left(\sum_j\sigma_jG^{\alpha_j}h_j\right)^n
 =\sum_{\abs{\bm m}=n}\frac{n!}{\bm m!}
 \bm\sigma^{\bm m}G^a h_{\bm m}.
\]
After division by \(G'=G\lambda\), the \(\bm m\)-term is
\(G^{-(1-a)}h_{\bm m}/\lambda\).  Apply
\eqref{eq:operator-factorization} successively with
\(k=1-a,2-a,\ldots,n-1-a\).  The final power is
\(G^{-(n-a)}\), and the factor \(n!\) cancels the one in the Lagrange
coefficient, leaving \eqref{eq:Vr-recurrence}--\eqref{eq:Umulti}.
\end{proof}

\begin{methodbox}[title={How to read the grading rule}]
Every perturbing scale \(G^{\alpha_j}h_j\) carries a positive defect
\(1-\alpha_j\) relative to \(G\).  Multiplication adds defects, so the inverse
sectors are indexed by the additive semigroup generated by
\(1-\alpha_1,\ldots,1-\alpha_s\).  If two multiindices have the same defect,
their coefficient blocks collide and must be added.  Logarithms, algebraic
powers, and Fourier phases remain in the coefficient algebra and do not alter
the exponential grade.
\end{methodbox}

\subsection{Differential coefficient algebras and Fourier closure}

A convenient formal setting is a differential algebra \(\Acal\) of ``slow''
coefficients closed under multiplication, \(D_x\), and multiplication by
\(\lambda^{-1}\).  The block operators \(\Lcal_k\) then preserve
\(\Acal\).  Typical choices include algebras generated by
\begin{equation}
 x^{-1},\quad (\log x)^{-1},\quad \log\log x,
 \quad \ee^{i\omega_1x},\ldots,\ee^{i\omega_rx},
 \label{eq:coefficient-generators}
\end{equation}
with formal power-logarithmic expansions in their amplitudes.  Products add
frequencies, so finite Fourier support remains finite in every fixed sector.

For the median subfactorial interpolation, take two coefficient modes
\begin{equation}
 h_+(x)=\tfrac12\ee^{i\pi x}J(x),
 \qquad
 h_-(x)=\tfrac12\ee^{-i\pi x}J(x),
 \qquad \alpha_+=\alpha_-=0.
 \label{eq:subfactorial-two-modes}
\end{equation}
A multiindex \((r,s)\) has grade \(r+s\) and frequency \((r-s)\pi\), which
immediately reproduces \eqref{eq:frequency-support}.

\subsection{Nested coefficient transseries}

Suppose a slow coefficient itself has a transseries
\begin{equation}
 h_j(x)\sim\sum_{\beta\in S_j}
 \zeta(x)^\beta h_{j,\beta}(x),
 \qquad \zeta(x)\to0.
 \label{eq:nested-coefficient-transseries}
\end{equation}
Substitution into \eqref{eq:Umulti} and collection of monomials enlarges the
sector semigroup by the grades carried by \(\zeta\).  Algebraically, the new
support is the Minkowski sum of the outer inversion semigroup and the internal
supports of the coefficients.  The Bell late-term effect in
\eqref{eq:Bell-sector-mixing} is precisely such a collision: the hidden
exponential scale in the asymptotic expansion of \(J\) enters at the same
\(y^{-2}\) grade as the second Lagrange block.

The practical conclusion is important.  A formally complete result should be
stated first in exact-coefficient form, such as \eqref{eq:Tn-y-form}; only then
should individual coefficient functions be asymptotically expanded, with a
truncation rule that respects all active scales.

\section{Validity, convergence, and remainder control}
\label{sec:validity}

The Lagrange formulas have two complementary interpretations.  They are
identities in a formal perturbation parameter, and they are analytic Taylor
series whenever the conjugated near-identity map is sufficiently small.  The
latter condition is automatically satisfied for the subfactorial at large
positive arguments.

\subsection{A local analytic criterion}

Let
\begin{equation}
 y=z+\varepsilon Q(z)
 \label{eq:analytic-near-identity}
\end{equation}
with \(Q\) analytic near a positive point \(y\).  The inverse value is the
fixed point of
\begin{equation}
 z=y-\varepsilon Q(z).
 \label{eq:fixed-point-z}
\end{equation}

\begin{proposition}[Convergence in the perturbation parameter]
\label{prop:analytic-convergence}
Suppose that for some disc \(B(y,R)\), constants \(M,\kappa>0\), and
\(\delta>0\),
\begin{equation}
 \sup_{B(y,R)}\abs{Q}\le M,
 \qquad
 \sup_{B(y,R)}\abs{Q'}\le\kappa,
 \qquad
 (1+\delta)M<R,
 \qquad
 (1+\delta)\kappa<1.
 \label{eq:contraction-hypotheses}
\end{equation}
Then for every \(\abs{\varepsilon}\le1+\delta\), equation
\eqref{eq:fixed-point-z} has a unique solution in \(B(y,R)\).  This solution is
analytic in \(\varepsilon\), and its Taylor series at \(\varepsilon=0\) is the
Lagrange--B\"urmann series.  In particular, the series converges at
\(\varepsilon=1\).
\end{proposition}
\begin{proof}
The right-hand side of \eqref{eq:fixed-point-z} maps the disc into itself and
has Lipschitz constant at most \((1+\delta)\kappa<1\).  The contraction mapping
theorem gives the unique fixed point.  Analytic dependence on \(\varepsilon\)
follows from the analytic implicit function theorem because
\(1+\varepsilon Q'(z)\ne0\).  The Taylor coefficients are those of the
Lagrange--B\"urmann formula; see, for example, \cite{Sokal2009}.
\end{proof}

\begin{proposition}[Application to the subfactorial]
\label{prop:subfactorial-convergence}
For the canonical interpolation \(\Dfrak=G+q\), the series
\eqref{eq:X-epsilon-series} converges at \(\varepsilon=1\) for every
sufficiently large positive \(y\).  Moreover, for each fixed \(N\), with
\(x=g(y)\),
\begin{equation}
 \Xfrak(y)=g(y)+\sum_{n=1}^{N}\Tcal_n(y)
 +\bigO_N\!\left(
 \frac{1}{y^{N+1}x^{N+1}\log x}
 \right).
 \label{eq:subfactorial-sector-remainder}
\end{equation}
\end{proposition}
\begin{proof}[Proof sketch]
For \(Q(z)=q(g(z))\), \Cref{lem:J-derivative-bound} and
\(G'(x)=G(x)\psi(x+1)\) give
\begin{equation}
 Q(y)=\bigO(x^{-1}),
 \qquad
 Q'(y)=\frac{q'(x)}{G'(x)}
 =\bigO\!\left(\frac{1}{yx\log x}\right).
 \label{eq:Q-smallness}
\end{equation}
The same estimates hold uniformly in a sufficiently small complex
neighborhood of \(y\), because \(g\), \(q\), and \(G\) are analytic there and
the neighborhood in the \(x\)-plane shrinks like
\(1/(yx^2\log x)\).  Thus \eqref{eq:contraction-hypotheses} holds for large
\(y\).  Repeated use of the operator formula
\eqref{eq:Tn-x-form}, together with fixed-order derivative bounds
\(q^{(r)}(x)=\bigO_r(x^{-1})\), gives
\(\Tcal_n=\bigO_n(y^{-n}x^{-n}/\log x)\).  Taylor's theorem in
\(\varepsilon\), with the same uniform bounds, yields
\eqref{eq:subfactorial-sector-remainder}.
\end{proof}

The expansion \eqref{eq:subfactorial-sector-remainder} is in the exponential
grade \(y^{-1}\), with exact coefficient functions.  If those functions are
subsequently replaced by divergent internal expansions, the nested warning of
\Cref{subsec:Bell-late-terms} applies.

\subsection{Residual certificates}

Formal asymptotics become rigorous numerical approximations once their
residual is checked.  This is particularly easy here because \(\Dfrak\) is
uniformly increasing on the principal branch.

\begin{proposition}[A posteriori error bound]
\label{prop:residual-certificate}
Let \(y\ge2\), and let \(\widehat{x}\ge3\).  Define the residual
\begin{equation}
 r=\Dfrak(\widehat{x})-y.
 \label{eq:residual}
\end{equation}
Then
\begin{equation}
 \abs{\widehat{x}-\Xfrak(y)}
 \le\frac{\abs r}{m_0},
 \qquad
 m_0=\frac6\ee\left(\frac{11}{6}-\gamma\right)
 -\left(\frac\pi4+\frac1{16}\right)>1.92.
 \label{eq:global-residual-bound}
\end{equation}
If a sharper lower bound \(m\le\Dfrak'(x)\) is known on the interval joining
\(\widehat{x}\) and \(\Xfrak(y)\), then \(m_0\) may be replaced by \(m\).
\end{proposition}
\begin{proof}
The proof of \Cref{thm:monotonicity} gives
\(\Dfrak'(x)\ge m_0\) for \(x\ge3\).  Apply the mean value theorem to
\(\Dfrak(\widehat{x})-\Dfrak(\Xfrak(y))=r\).
\end{proof}

For large \(y\), a local derivative estimate is much stronger:
\begin{equation}
 \Dfrak'(x)=y\log x\,(1+o(1)),
 \qquad x=\Xfrak(y).
 \label{eq:Dprime-asymptotic}
\end{equation}
Thus the error is approximately the residual divided by \(y\log x\).  One
Newton step
\begin{equation}
 \widehat{x}_{\mathrm{new}}
 =\widehat{x}-\frac{\Dfrak(\widehat{x})-y}{\Dfrak'(\widehat{x})}
 \label{eq:Newton-D}
\end{equation}
usually converts a short transseries truncation into a high-precision value.

\subsection{Balanced truncation}

There are two independent truncation indices:
\begin{enumerate}
\item \(M\), the number of zero-sector inverse-gamma blocks in
\eqref{eq:g-outer-main};
\item \(N\), the number of exponential Lagrange sectors in
\eqref{eq:subfactorial-sector-remainder}.
\end{enumerate}
The zero-sector error should be made no larger than the last retained
exponential sector.  If an approximate value \(\widetilde g\) is used in the
phase \(\cos(\pi g)\), then its absolute error propagates directly into the
phase.  Consequently, retaining an extremely small \(y^{-N}\) correction is
pointless unless \(g\) has been computed to comparable absolute accuracy.  In
practice, one either takes enough terms in \eqref{eq:g-outer-main} or performs
one Newton correction for \(G(x)=y\) before evaluating the exponential
sectors.

\section{Numerical validation}
\label{sec:numerics}

The following computations used high-precision arithmetic.  To isolate the
exponential hierarchy, the exact numerical inverse \(g(y)\) of \(G\) was used
inside \(\Tcal_1,\Tcal_2,\Tcal_3\).  The final column separately reports the
error of the four-block outer approximation
\(X-\tfrac12+\sum_{m=1}^4a_m/X^{2m-1}\).

\begin{table}[htbp]
\centering
\small
\caption{Successive exponential sectors for generic arguments.}
\label{tab:generic-numerics}
\begin{tabular}{@{}rllll@{}}
\toprule
\(y\) & \(\abs{g-\Xfrak}\) &
\(\abs{g+\Tcal_1-\Xfrak}\) &
\(\abs{g+\Tcal_1+\Tcal_2-\Xfrak}\) &
outer \(g\) error \\
\midrule
\(10^2\)  & \(7.6371\times10^{-5}\)  & \(1.8806\times10^{-7}\) &
\(4.7609\times10^{-10}\) & \(5.4645\times10^{-11}\) \\
\(10^6\)  & \(3.3545\times10^{-8}\)  & \(4.6652\times10^{-18}\) &
\(2.2653\times10^{-22}\) & \(2.8457\times10^{-13}\) \\
\(10^{20}\) & \(1.8356\times10^{-23}\) & \(7.3241\times10^{-45}\) &
\(2.6668\times10^{-66}\) & \(2.4032\times10^{-16}\) \\
\bottomrule
\end{tabular}
\end{table}

The corresponding exact inverse values are
\begin{align*}
 \Xfrak(10^2)&=5.4680181917506673490048452082893058\ldots,\\
 \Xfrak(10^6)&=9.8768464792043069546910118994035847\ldots,\\
 \Xfrak(10^{20})&=21.542599928440196883351107787710195\ldots.
\end{align*}
Adding \(\Tcal_3\) reduces the three errors in the third data column further
to approximately \(1.23\times10^{-12}\), \(6.08\times10^{-30}\), and
\(4.02\times10^{-80}\), respectively.  The table also illustrates balanced
truncation: for very large \(y\), a short zero-sector approximation can be the
dominant numerical error even after the exponential sectors have become
extraordinarily accurate.

At the exact derangement values \(y=D_n\), the continuous inverse must return
\(n\).  The cancellation of the inverse-gamma displacement is shown in
\Cref{tab:integer-numerics}.

\begin{table}[htbp]
\centering
\small
\caption{Recovery of integer arguments at \(y=D_n\).  Entries are signed
errors relative to \(n\).}
\label{tab:integer-numerics}
\begin{tabular}{@{}r r r r r@{}}
\toprule
\(n\) & \(D_n\) & \(g(D_n)-n\) & \(g+\Tcal_1-n\) &
\(g+\Tcal_1+\Tcal_2-n\) \\
\midrule
5 & 44 & \(-1.93565\times10^{-3}\) & \(3.91873\times10^{-6}\) &
\(2.51266\times10^{-8}\) \\
10 & 1334961 & \(2.67167\times10^{-8}\) & \(9.13999\times10^{-16}\) &
\(-5.19670\times10^{-23}\) \\
20 & 895014631192902121 & \(1.68472\times10^{-20}\) &
\(4.43897\times10^{-40}\) & \(-7.92300\times10^{-60}\) \\
\bottomrule
\end{tabular}
\end{table}

The leading integer displacement follows directly from the first sector.
Because \(D_n=G(n)+(-1)^nJ(n)\), inversion of \(G\) alone gives
\begin{equation}
 g(D_n)-n
 =\frac{(-1)^nJ(n)}{G(n)\psi(n+1)}
 +\bigO(G(n)^{-2})
 \sim \frac{(-1)^n\ee}{n!\,\psi(n+1)}
 \left(\frac1n-\frac2{n^2}+\cdots\right).
 \label{eq:integer-displacement}
\end{equation}
The term \(\Tcal_1(D_n)\) cancels this displacement to the next exponential
grade.

\subsection{A check on optimal Bell truncation}

For orientation, truncating \eqref{eq:J-Bell-explicit} near
\(M\approx x\log x\) gives the errors shown below.  The last column divides
the error by the scale \(\sqrt{x/\log x}/G(x)\) predicted by
\eqref{eq:Bell-least-term-scale}.

\begin{table}[htbp]
\centering
\small
\caption{Late-term scale of the Bell expansion of \(J\).}
\label{tab:Bell-late-numerics}
\begin{tabular}{@{}r r r r@{}}
\toprule
\(x\) & \(M\) & Bell truncation error & normalized error \\
\midrule
5  & 8   & \(5.25774\times10^{-3}\)  & 0.1317 \\
10 & 23  & \(2.20217\times10^{-7}\)  & 0.1411 \\
20 & 60  & \(4.23592\times10^{-19}\) & 0.1467 \\
40 & 148 & \(1.65201\times10^{-48}\) & 0.1506 \\
\bottomrule
\end{tabular}
\end{table}

The stable normalized values support the sector-mixing estimate
\eqref{eq:Bell-sector-mixing}.  For numerical work, however, the convergent
inverse-factorial series \eqref{eq:J-inverse-factorial} removes this issue
entirely.

\section{Symbolic and numerical construction}
\label{sec:algorithms}

The formulas above support several complementary workflows.  The choice
depends on whether the goal is a formal transseries, a fast numerical value, or
a certified integer inverse.

\subsection{A robust numerical algorithm}

\begin{algorithm}[Inverse subfactorial for large \(y\)]
\label{alg:numerical-inverse}
Given \(y\ge2\), an outer order \(M\), and an exponential-sector order
\(N\):
\begin{enumerate}
\item Compute
\[
 A=\log\!\frac{\ee y}{\sqrt{2\pi}},\qquad
 w=\W_0(A/\ee),\qquad X=A/w,
\]
and form the \(M\)-block approximation to \(g(y)\) from
\eqref{eq:g-outer-main}.
\item If more accuracy is needed, apply one or two Newton steps to
\(G(x)-y=0\):
\[
 x\leftarrow x-\frac{G(x)-y}{G(x)\psi(x+1)}.
\]
The resulting value is the numerical base point \(x\approx g(y)\).
\item Evaluate \(J(x)\) by the convergent recurrence
\eqref{eq:J-term-recurrence}.  Evaluate \(J'\), \(J''\), and higher
derivatives either by differentiating the integral or by automatic
differentiation of the same convergent series.
\item Form \(q,q',q''\) from \eqref{eq:q-cJ}--\eqref{eq:qpp-cJ}, and evaluate
\(\Tcal_1,\ldots,\Tcal_N\) from \eqref{eq:Tn-x-form}.  For \(N\le3\), use
\eqref{eq:T1-polygamma}--\eqref{eq:T3-polygamma} directly.
\item Return
\[
 \widehat{x}=x+\sum_{n=1}^{N}\Tcal_n(y).
\]
Optionally apply one Newton step \eqref{eq:Newton-D} to the full interpolation.
\item Compute the residual \(r=\Dfrak(\widehat{x})-y\).  Use
\eqref{eq:global-residual-bound}, or a sharper local derivative bound, as a
certificate.
\end{enumerate}
\end{algorithm}

For the integer generalized inverse \(N(y)\), it is enough to certify that the
continuous error is smaller than the distance from \(\widehat{x}\) to the
nearest integer boundary.  Then \(N(y)=\lfloor\widehat{x}\rfloor\) by
\Cref{cor:integer-inverse}.

\subsection{Direct symbolic generation of the exponential blocks}

The operator formula \eqref{eq:Tn-x-form} is ideal for a computer algebra
system.  The following pseudocode keeps \(G\) and \(q\) symbolic and generates
arbitrary sectors.

\begin{lstlisting}
# Symbolic pseudocode for T_n
# Dx denotes differentiation with respect to x.

def sector(n, G, q, x):
    A = diff(G, x)
    expr = q**n / A
    for k in range(n - 1):
        expr = diff(expr, x) / A
    return (-1)**n * expr / factorial(n)

# Substitute x = g(y) and G(g(y)) = y afterwards.
\end{lstlisting}

A useful simplification strategy is:
\begin{enumerate}
\item replace \(G^{(r)}\) by complete Bell polynomials in the derivatives of
\(\log G\);
\item for the gamma dominant term, replace those derivatives by polygamma
functions;
\item factor the resulting power of \(G(x)=y\);
\item only then expand \(J\) and the polygamma functions asymptotically.
\end{enumerate}
This ordering preserves the exponential grading and avoids premature
expansion of the Bell coefficient.

\subsection{Symbolic generation of the zero-sector coefficients}

The triangular identity \eqref{eq:eta-master} is similarly easy to automate.
A concise implementation is shown below.

\begin{lstlisting}
# SymPy-style pseudocode for a_m(L)
# kappa(r) is the centered Stirling coefficient.

t, L = symbols('t L')
eta = 0
coefficients = []

for m in range(1, M + 1):
    a = symbol(f'a{m}')
    trial = eta + a*t**(2*m)
    equation = (trial*(L - 1)
                + (1 + trial)*log(1 + trial))/t
    equation += sum(kappa(r)*t**(2*r - 1)
                    *(1 + trial)**(-(2*r - 1))
                    for r in range(1, m + 1))
    c = coefficient(series(equation, t, 0, 2*m),
                    t, 2*m - 1)
    a_value = solve(c == 0, a)[0]
    coefficients.append(simplify(a_value))
    eta += a_value*t**(2*m)
\end{lstlisting}

Because the equation at each stage is linear in the new coefficient, this
algorithm is stable as a formal calculation and scales substantially better
than a blind reversion of the original gamma equation.

\subsection{A triangular shift recurrence}

There is an alternative to Lagrange--B\"urmann inversion.  Set
\begin{equation}
 x=g(y)+\Delta(y)
 \label{eq:Delta-shift}
\end{equation}
and expand
\begin{equation}
 0=G(g+\Delta)-G(g)+H(g+\Delta).
 \label{eq:Delta-equation}
\end{equation}
If \(\Delta=\sum_{n\ge1}y^{-n}U_n(g)\), then comparison of powers of \(y^{-1}\)
produces a triangular recurrence for \(U_n\).  The first equations are
\begin{align}
 0={}&\lambda U_1+H,
 \label{eq:shift-U1}\\
 0={}&\lambda U_2+
rac12(\lambda^2+\lambda')U_1^2+H'U_1,
 \label{eq:shift-U2}
\end{align}
where all functions are evaluated at \(x=g(y)\).  Solving gives
\(U_1=-H/\lambda\) and the same \(U_2\) as \eqref{eq:U2}.  This approach is
useful for hand calculations and for proving uniqueness of the coefficient
blocks.  The operator formula is usually more compact at high order.

\section{Examples and extensions of the general theory}
\label{sec:extensions}

\subsection{Factorial-type dominant maps}

Suppose
\begin{equation}
 \log G(x)\sim x\log x-x+\beta\log x+c
 +\sum_{k\ge1}\frac{c_k}{x^k}.
 \label{eq:factorial-type-G}
\end{equation}
Then the leading inverse is Lambert-\(W\)-controlled.  Centering the variable
to eliminate as many even or odd correction powers as possible produces a
triangular zero-sector recurrence exactly like \eqref{eq:eta-master}.  Once
\(g\) is known, any additive tail
\begin{equation}
 H(x)\sim\sum_{\omega\in\Omega}\ee^{i\omega x}
 x^{\rho_\omega}(\log x)^{\tau_\omega}
 \sum_{k\ge0}\frac{h_{\omega,k}}{x^k}
 \label{eq:factorial-additive-tail}
\end{equation}
is inverted by \eqref{eq:Un-formula}.  The \(n\)-th inverse sector has
exponential grade \(G(g(y))^{-n}=y^{-n}\), finite frequency support in the
\(n\)-fold sumset of \(\Omega\), and coefficient asymptotics generated by
\(\lambda^{-1}D_x-k\).

The subfactorial is the model case with \(\beta=1/2\),
\(\Omega=\{\pi,-\pi\}\), and \(H=q\sim x^{-1}\cos(\pi x)\).

\subsection{Perturbations of fractional dominant size}

Let
\begin{equation}
 F(x)=G(x)+G(x)^\alpha h(x),\qquad \alpha<1.
 \label{eq:fractional-perturbation}
\end{equation}
The first inverse correction is
\begin{equation}
 -\frac{G(x)^\alpha h(x)}{G'(x)}
 =-G(x)^{-(1-\alpha)}\frac{h(x)}{\lambda(x)}.
 \label{eq:fractional-first}
\end{equation}
Thus the natural small parameter is \(y^{-(1-\alpha)}\), not \(y^{-1}\).
The \(n\)-th pure sector has grade \(n(1-\alpha)\), while mixed perturbations
obey the additive rule \eqref{eq:sector-grade}.  Irrational defects produce a
noninteger but still well-ordered exponential scale.

\subsection{Multiplicative perturbations}

A relatively small multiplicative correction
\begin{equation}
 F(x)=G(x)(1+r(x)),\qquad r(x)\to0,
 \label{eq:multiplicative-perturbation}
\end{equation}
formally corresponds to \(\alpha=1\), so its grade is not supplied by a power
of \(y^{-1}\).  Instead, the internal scale of \(r\) supplies the ordering.
The first correction is
\begin{equation}
 F^{-1}(y)-g(y)\sim-\frac{r(g(y))}{\lambda(g(y))}.
 \label{eq:multiplicative-first}
\end{equation}
This case is covered by the exact Lagrange formula and by the nested
coefficient formalism of \eqref{eq:nested-coefficient-transseries}.  It shows
why a general transseries theory must track both powers of the dominant
function and independent small coefficient scales.

\subsection{Several saddles or lateral continuations}

Suppose an interpolation has lateral corrections
\begin{equation}
 H(x)=\sum_{j=1}^{s}S_j\ee^{-\Phi_j(x)}a_j(x),
 \label{eq:several-saddles}
\end{equation}
where \(S_j\) are Stokes or interpolation parameters.  If
\(\ee^{-\Phi_j(x)}=G(x)^{-\mu_j}\zeta_j(x)\) with slow \(\zeta_j\), then
\(\alpha_j=1-\mu_j\), and \Cref{thm:multiscale} gives the complete mixed
sector lattice.  Products generate actions
\(\sum_jm_j\Phi_j\), while the block operators generate their prefactors.
This is the inversion analogue of the usual additive action semigroup in
transseries and resurgence; compare the general framework in
\cite{vdHoeven2006}.

\subsection{Interpolation dependence revisited}

Let two admissible interpolations be
\begin{equation}
 F_1(x)=G(x)+H_1(x),\qquad
 F_2(x)=G(x)+H_2(x),
 \label{eq:two-interpolations}
\end{equation}
with \(H_1(n)=H_2(n)\) at all integers and both \(H_j=\bigO(x^{-1})\).  Their
inverse zero sectors agree to every algebraic order in \(X^{-1}\), because both
are anchored by the same \(G\).  Their first difference is
\begin{equation}
 F_1^{-1}(y)-F_2^{-1}(y)
 =-\frac{H_1(g(y))-H_2(g(y))}{y\lambda(g(y))}
 +\bigO(y^{-2}).
 \label{eq:interpolation-inverse-difference}
\end{equation}
Thus interpolation ambiguity is itself an exponentially small transseries
parameter.  The canonical median choice \eqref{eq:median-D} fixes that
parameter by a transparent contour prescription.

\section{Interpretation of the result}
\label{sec:interpretation}

\subsection{Why the inverse is much closer to inverse gamma than the function is to gamma}

The additive discrepancy \(q(x)\) is only algebraically small,
\(q(x)=\bigO(x^{-1})\), whereas the derivative of the dominant term is
exponentially large:
\begin{equation}
 G'(x)=G(x)\psi(x+1)\asymp G(x)\log x.
 \label{eq:large-derivative}
\end{equation}
Therefore the inverse displacement is the quotient
\begin{equation}
 \Xfrak(y)-g(y)
 \sim-\frac{q(g(y))}{G'(g(y))}
 =\bigO\!\left(\frac1{y g(y)\log g(y)}\right).
 \label{eq:inverse-displacement-scale}
\end{equation}
This elementary derivative mechanism explains the enormous accuracy of
inverse gamma as an approximation to inverse subfactorial.

\subsection{Why the cosine phase cannot be replaced by \texorpdfstring{\((-1)^x\)}{(-1) raised to x}}

At integers, \(\cos(\pi n)=(-1)^n\).  Away from the integers, however, the
real phase is an essential part of the interpolation.  The notation
\((-1)^x\) is branch-dependent and generally complex.  The two lateral
continuations naturally carry \(\ee^{\pm i\pi x}\), while their real median
carries \(\cos(\pi x)\).  The finite Fourier structure of all higher sectors
is inherited from this choice.

\subsection{Universality and nonuniversality}

The result separates naturally into two parts.
\begin{itemize}
\item \emph{Universal part:} the Lambert-anchored inverse-gamma zero sector,
the Lagrange--B\"urmann operator calculus, the sector grading rule, and the
residual certificates.
\item \emph{Interpolation-specific part:} the precise phase and coefficient
\(q(x)=\cos(\pi x)J(x)\), hence the numerical values of the exponentially small
sectors.
\end{itemize}
For any other interpolation whose additive discrepancy from \(G\) is
\(\bigO(x^{-1})\), the same zero sector applies and the same operator formulas
compute the modified sectors after replacing \(q\).

\subsection{What is genuinely transserial here}

A finite inverse-gamma expansion alone is an ordinary asymptotic series.  The
full object contains a strict hierarchy
\begin{equation}
 X\gg1\gg X^{-1}\gg X^{-2}\gg\cdots
 \gg y^{-1}x^{-1}\gg y^{-2}x^{-2}\gg\cdots,
 \label{eq:hierarchy-display}
\end{equation}
with logarithmic coefficients and oscillatory monomials.  The hierarchy is
closed under multiplication, differentiation, and composition, and its
support is locally finite at each exponential grade.  These are the defining
operational features of a transseries, even though the present positive-real
problem requires no Stokes switching once the median interpolation has been
fixed.

\section{Conclusions}
\label{sec:conclusion}

The inverse subfactorial problem has a clean answer only after the continuous
object to be inverted is specified.  The median incomplete-gamma interpolation
\[
 \Dfrak(x)=\Gamma(x+1)/\ee+\cos(\pi x)J(x)
\]
is natural, real, exact at every integer, and monotone on the large branch.
Its Bell coefficient \(J\) admits both a convergent inverse-factorial series
and a real-variable Bell expansion with a signed, explicit remainder.  This
makes it possible to separate numerical evaluation from asymptotic analysis.

The dominant inverse is controlled by a Lambert-\(W\) anchor.  Centering at
\(x+1/2\) yields an unusually sparse odd-power expansion, and the coefficient
blocks satisfy a simple triangular recurrence.  The derangement correction is
beyond all orders of that zero sector.  Conjugation by the dominant gamma map
turns the remaining problem into a near-identity inversion, and the
Lagrange--B\"urmann formula gives every exponential sector explicitly.

The reusable theory is summarized by two formulas.  For a single additive
perturbation,
\[
 U_n=\frac{(-1)^n}{n!}
 (\lambda^{-1}D_x-(n-1))\cdots(\lambda^{-1}D_x-1)
 \left(\frac{H^n}{\lambda}\right),
 \qquad \lambda=D_x\log G.
\]
For several relative scales \(G^{\alpha_j}h_j\), a multiindex \(\bm m\) has
inverse-sector grade
\[
 \nu(\bm m)=\sum_jm_j(1-\alpha_j).
\]
These formulas apply to a broad range of factorial-type, gamma-type, and
multi-saddle transseries.

Finally, the Bell coefficient has its own divergent power expansion.  Its late
terms reach the next exponential grade, so a completely expanded inverse is a
nested transseries.  Keeping \(J\) exact, or evaluating it through its
convergent inverse-factorial series, preserves a clean sector decomposition and
provides a reliable computational route.  This distinction between
exact-coefficient sectors and internally expanded sectors is likely the most
portable lesson of the subfactorial example.

\appendix

\section{Derivation of the centered Stirling coefficients}
\label{app:centered-Stirling}

The shifted Stirling expansion may be written
\begin{equation}
 \log\Gamma(u+h)
 \sim u\log u-u+(h-\tfrac12)\log u+\tfrac12\log(2\pi)
 +\sum_{n\ge1}\frac{(-1)^{n+1}\Bell^{\mathrm{Ber}}_{n+1}(h)}
 {n(n+1)u^n},
 \label{eq:shifted-Stirling-app}
\end{equation}
where \(\Bell^{\mathrm{Ber}}_{n}(h)\) is the Bernoulli polynomial.  Set
\(h=1/2\).  Odd Bernoulli polynomials vanish at \(1/2\), and
\begin{equation}
 \Bell^{\mathrm{Ber}}_{2m}(1/2)
 =(2^{1-2m}-1)\Bell^{\mathrm{Ber}}_{2m}.
 \label{eq:Bernoulli-half}
\end{equation}
Since \(u+h=x+1\), equations
\eqref{eq:shifted-Stirling-app}--\eqref{eq:Bernoulli-half} reduce to
\eqref{eq:centered-Stirling}--\eqref{eq:kappa-def}.

The first few centered terms are
\begin{equation}
 \log\Gamma(x+1)
 \sim u(\log u-1)+\frac12\log(2\pi)
 -\frac1{24u}+\frac7{2880u^3}-\frac{31}{40320u^5}
 +\frac{127}{215040u^7}-\frac{511}{608256u^9}+\cdots.
 \label{eq:centered-Stirling-explicit-app}
\end{equation}

\section{Coefficient extraction for the outer inverse}
\label{app:outer-coefficients}

Let
\begin{equation}
 \mathscr{E}(t,\eta)=
 t^{-1}\{\eta(L-1)+(1+\eta)\log(1+\eta)\}
 +\sum_{r\ge1}\kappa_rt^{2r-1}(1+\eta)^{-(2r-1)}.
 \label{eq:E-functional}
\end{equation}
The coefficient recurrence can be expressed without introducing the temporary
symbol \(\rho_m\):
\begin{equation}
 a_m(L)=-\frac1L
 [t^{2m-1}]\,
 \mathscr{E}\!\left(t,\sum_{j=1}^{m-1}a_j(L)t^{2j}\right).
 \label{eq:am-coefficient-extraction}
\end{equation}
Here \([t^k]\) denotes coefficient extraction.  Formula
\eqref{eq:am-coefficient-extraction} is valid because the omitted contribution
from \(a_mt^{2m}\) is exactly \(La_mt^{2m-1}\).

For example, at first order,
\[
 0=La_1+\kappa_1,
 \qquad a_1=-\kappa_1/L=1/(24L).
\]
At second order,
\[
 0=La_2+\frac12a_1^2-\kappa_1a_1+\kappa_2,
\]
which simplifies to \eqref{eq:a2}.  Higher orders are no harder conceptually;
only the polynomial bookkeeping grows.

\section{Low-order universal blocks in expanded form}
\label{app:universal-blocks}

For reference, let \(D=D_x\).  Expanding \eqref{eq:U2} gives
\begin{equation}
 U_2=
rac{HH'}{\lambda^2}
 -\frac{H^2}{2\lambda}
 -\frac{H^2\lambda'}{2\lambda^3}.
 \label{eq:U2-expanded}
\end{equation}
This agrees with \eqref{eq:T2-general} after using
\(G''/G=\lambda^2+\lambda'\).

The third block can be left compactly as \eqref{eq:U3}; an expanded form is
\begin{align}
 U_3={}&-\frac{H(H')^2}{\lambda^3}
 -\frac{H^2H''}{2\lambda^3}
 +\frac{3H^2H'}{2\lambda^2}
 +\frac{3H^2H'\lambda'}{2\lambda^4}
 \notag\\
 &-\frac{H^3}{3\lambda}
 -\frac{H^3\lambda'}{2\lambda^3}
 -\frac{H^3(\lambda')^2}{2\lambda^5}
 +\frac{H^3\lambda''}{6\lambda^4}.
 \label{eq:U3-expanded}
\end{align}
Substitution \(H=q\), \(\lambda=\psi(x+1)\) and multiplication by \(y^{-3}\)
recovers \eqref{eq:T3-polygamma} after elementary simplification.

\section{Bell and Fourier coefficient generation}
\label{app:Bell-Fourier}

The power coefficients of \(J\) are
\begin{equation}
 j_k=(-1)^{k-1}\Bell_k,
 \qquad
 J(x)\sim\sum_{k\ge1}j_kx^{-k}.
 \label{eq:jk-def}
\end{equation}
For a lateral mode
\begin{equation}
 q_\sigma(x)=\frac12\ee^{i\sigma\pi x}J(x),
 \qquad \sigma\in\{+1,-1\},
 \label{eq:q-sigma}
\end{equation}
one has
\begin{equation}
 D_x^r q_\sigma(x)
 \sim\frac12\ee^{i\sigma\pi x}
 \sum_{s=0}^{r}\binom{r}{s}(i\sigma\pi)^{r-s}
 (-1)^s\sum_{k\ge1}
 \frac{(k)_s j_k}{x^{k+s}}.
 \label{eq:q-sigma-derivatives}
\end{equation}
This formula, together with the Bernoulli expansions of the polygamma
functions, gives a direct coefficient-level algorithm for any fixed Fourier
mode in any fixed exponential sector.

For the median mode \(q=q_++q_-\), products are collected by the binomial
identity
\begin{equation}
 q(x)^n=2^{-n}J(x)^n
 \sum_{r=0}^{n}\binom{n}{r}\ee^{i(n-2r)\pi x}.
 \label{eq:q-power-Fourier}
\end{equation}
This makes the support \eqref{eq:frequency-support} explicit before any
differentiation is performed.

\section{Notation and hierarchy at a glance}
\label{app:notation}

\begin{longtable}{@{}p{0.20\textwidth}p{0.72\textwidth}@{}}
\toprule
Symbol & Meaning \\
\midrule
\endfirsthead
\toprule
Symbol & Meaning \\
\midrule
\endhead
\(D_n={!n}\) & Number of derangements of \(n\) objects. \\
\(G(x)\) & Dominant interpolation \(\Gamma(x+1)/\ee\). \\
\(J(x)\) & Positive Bell-tail coefficient
\(\int_0^1\ee^{-t}(1-t)^x\,\dd t\). \\
\(q(x)\) & Oscillatory additive correction \(\cos(\pi x)J(x)\). \\
\(\Dfrak(x)\) & Canonical real interpolation \(G(x)+q(x)\). \\
\(g(y)\) & Increasing large inverse of \(G\). \\
\(\Xfrak(y)\) & Principal inverse of \(\Dfrak\) on \([3,\infty)\). \\
\(A,w,X,L\) & Lambert anchor variables from \eqref{eq:AXwL}. \\
\(p,p_1,p_2\) & Digamma and first two higher polygamma functions at \(x+1\). \\
\(\lambda\) & Logarithmic derivative \(G'/G\); for gamma, \(\lambda=\psi(x+1)\). \\
\(\Tcal_n\) & \(n\)-th exponential inverse sector, carrying \(y^{-n}\). \\
\(\Lcal_k\) & Universal block operator \(\lambda^{-1}D_x-k\). \\
\(\Bell_k\) & Bell number; Bernoulli numbers are written
\(\Bell_k^{\mathrm{Ber}}\). \\
\bottomrule
\end{longtable}

The main hierarchy is
\begin{equation}
 g(y)=X-\frac12+\bigO(X^{-1}L^{-1}),
 \qquad
 \Tcal_n(y)=\bigO_n\!\left(
 y^{-n}g(y)^{-n}(\log g(y))^{-1}
 \right).
 \label{eq:hierarchy-summary}
\end{equation}
Thus the zero sector is expanded in odd inverse powers of \(X\), while the
nonzero sectors are graded by powers of \(y^{-1}\) and carry internal
power-logarithmic and Fourier expansions.

\begin{thebibliography}{99}

\bibitem{Hassani2020}
M.~Hassani,
``Derangements and Alternating Sum of Permutations by Integration,''
\emph{Journal of Integer Sequences} \textbf{23} (2020), no.~7,
Article 20.7.8, 9 pp.
\url{https://cs.uwaterloo.ca/journals/JIS/VOL23/Hassani/hassani5.html}

\bibitem{DLMF4}
NIST Digital Library of Mathematical Functions,
\S4.13, ``Lambert \(W\)-Function,'' Release 1.2.7 (2026).
\url{https://dlmf.nist.gov/4.13}

\bibitem{DLMF5}
NIST Digital Library of Mathematical Functions,
\S5.11, ``Asymptotic Expansions'' for the gamma function,
Release 1.2.7 (2026).
\url{https://dlmf.nist.gov/5.11}

\bibitem{DLMF26}
NIST Digital Library of Mathematical Functions,
\S26.7, ``Set Partitions: Bell Numbers,'' Release 1.2.7 (2026).
\url{https://dlmf.nist.gov/26.7}

\bibitem{Corless1996}
R.~M.~Corless, G.~H.~Gonnet, D.~E.~G.~Hare, D.~J.~Jeffrey, and D.~E.~Knuth,
``On the Lambert \(W\) Function,''
\emph{Advances in Computational Mathematics} \textbf{5} (1996), 329--359.
\url{https://doi.org/10.1007/BF02124750}

\bibitem{Sokal2009}
A.~D.~Sokal,
``A Ridiculously Simple and Explicit Implicit Function Theorem,''
\emph{S\'eminaire Lotharingien de Combinatoire} \textbf{61A}
(2009/2011), Article B61Ad.
\url{https://arxiv.org/abs/0902.0069}

\bibitem{MoserWyman1955}
L.~Moser and M.~Wyman,
``An Asymptotic Formula for the Bell Numbers,''
\emph{Transactions of the Royal Society of Canada}, Series III,
\textbf{49} (1955), Section III, 49--54.

\bibitem{deBruijn1961}
N.~G.~de Bruijn,
\emph{Asymptotic Methods in Analysis},
North-Holland, Amsterdam, 1961.

\bibitem{Olver1997}
F.~W.~J.~Olver,
\emph{Asymptotics and Special Functions},
A.~K.~Peters, Wellesley, MA, 1997; originally published by Academic Press,
1974.

\bibitem{Comtet1974}
L.~Comtet,
\emph{Advanced Combinatorics},
D.~Reidel, Dordrecht, 1974.

\bibitem{FlajoletSedgewick2009}
P.~Flajolet and R.~Sedgewick,
\emph{Analytic Combinatorics},
Cambridge University Press, Cambridge, 2009.

\bibitem{vdHoeven2006}
J.~van der Hoeven,
\emph{Transseries and Real Differential Algebra},
Lecture Notes in Mathematics 1888, Springer, Berlin, 2006.
\url{https://doi.org/10.1007/3-540-35590-1}

\end{thebibliography}

\end{document}
